\begin{figure}
\begin{lstlisting}[
    basicstyle=\footnotesize\ttfamily,
    xleftmargin=2.5em,
    numbers=left]
# Extracts a set of $n$ learned abstractions.
# $E$ is an e-graph containing the set of programs with
# all possible function definitions and applications
# $r$ is the e-class id of the root node of the e-graph
# $n$ is number of libraries to learn at a time
def choose_abstraction($E, r, n$):
    # $CS$ maps e-classes IDs to CostSets
    $CS$ = $\{\}$
    # run till fixpoint
    for $C$ in $E$.e-classes:
        for $N$ in $C$.e-nodes:
            $CS[N]$ = cs-prune(cs-unify(cs-union($CS[N]$, cs-node($N$))))
    best = $min_{full-cost}(CS[r])$
    return best.libs

\end{lstlisting}
\caption{The extraction algorithm of \babble.}
\label{fig:algo-extract}
\end{figure}

\begin{figure}
  \centering
  \begin{subfigure}{\textwidth}
\begin{lstlisting}[
    basicstyle=\footnotesize\ttfamily,
    xleftmargin=2.5em,
    numbers=left]
# The cost of an expression, e, can be expressed by a LibrarySelection,
# which is the set of libraries used in e, along with their associated costs.
# e's cost is then its AST size, not including the size of the libraries used within it.
# We define an ordering between library selections:
# comparing library selection $a$ and $b$ is defined as first comparing
# $a$.cost to $b$.cost, then comparing full_cost($a$) to full_cost(b).
datatype LibrarySelection:
    cost: int
    libs: (abstraction * cost) set

# This function computes the cost of the expression a library selection corresponds to,
# where abstraction definitions are only counted once.
# $l$ is a library selection.
def full_cost($l$):
    return $l$.cost + sum([$c$ | _, $c$ in $l$.libs])

# A CostSet is a set of possible library selections that can be used at a particular e-class.
# The set must be sorted in ascending order.
datatype CostSet:
    cost_set: LibrarySelection set
\end{lstlisting}
\end{subfigure}

\begin{subfigure}{\textwidth}
\begin{lstlisting}[
    basicstyle=\footnotesize\ttfamily,
    xleftmargin=2.5em,
    numbers=left]
# Performs partial order reduction over a CostSet, $a$.
def cs-unify($a$):
    i = 0
    while i < len($a$):
        j = i + 1
        while j < len($a$):
            if $a$[i].libs is subset of $a$[j].libs:
                remove $a$[j]
            else:
                j += 1
        i += 1
    return $a$
\end{lstlisting}
\end{subfigure}

\begin{subfigure}{\textwidth}
\begin{lstlisting}[
    basicstyle=\footnotesize\ttfamily,
    xleftmargin=2.5em,
    numbers=left]
# Performs a product between two cost sets, such that paired library selections
# have their libraries combined (via union) and their expression costs added.
# This operation is performed when an e-node has multiple children e-classes;
# the component libraries from each e-class
# must be combined via union, and their expression costs must be included.
# $a$ and $b$ are two input CostSets
  def cs-product($a, b$):
    result = $\{\}$
    for a-sel in $a$:
        for b-sel in $b$:
            # the union preserves lowest cost of abstraction
            new-libs = a-sel.libs $\cup$ b-sel.libs
            new-cost = a-sel.cost + b-sel.cost
            results.insert(new-libs, new-cost)
    return cs-unify(result)

\end{lstlisting}
\end{subfigure}
\caption{Helper functions for extraction in \babble.}
\label{fig:helper1}
\end{figure}

\begin{figure}
\begin{lstlisting}[
    basicstyle=\footnotesize\ttfamily,
    xleftmargin=2.5em,
    numbers=left]
# A variant of cost set products used when encountering an e-node that defines
# an abstraction. In this case, a product must be performed which adds an
# abstraction to all library selections in the resulting cost set.
# a is the CostSet of the abstraction definition e-class.
# e is an identifier for the abstraction.
# b is the CostSet of the rest of the expression.
def cs-add-lib($a, e, b$):
    result = $\{\}$
    for a-sel in $a$:
        for b-sel in $b$:
            # the union preserves lowest cost of abstraction
            a-size = min([ls-full-cost(l) | l in a-sel])
            new-libs = union(a-sel.libs, b-sel.libs, [($e$, a-size)])
            new-cost = b-sel.cost
            result.insert(new-libs, new-cost)
    # returns the modified cost set product of a and b
    return cs-unify(result)

# Performs a union between two CostSets while maintaining the sorted invariant.
# $a$ and $b$ are CostSet
def cs-union($a, b$):
    result = $a$
    ix = 0
    for b-elem in $b$:
        while ix < len($a$) and elem >= $a$[ix]:
            ix += 1
        result.insert(b-elem)
        ix += 1
    return result
\end{lstlisting}
\caption{Helper functions for extraction in \babble.}
\label{fig:helper2}
\end{figure}

\begin{figure}
\begin{lstlisting}[
    basicstyle=\footnotesize\ttfamily,
    xleftmargin=2.5em,
    numbers=left]
// TODO: Not sure I follow this.
# $s$ is the input e-node and $CS$ the current map from e-class ids to CostSet.
def cs-node($s, CS$):
    ($a_0, ..., a_i$) = children($s$)
    if len(children(s)) == 0:
        return [($\{\}$, 1)]
    # if $s$ introduces a library function
    elif form($s$).matches()
    let f = \x. abs in body

        if there exists a non-empty cost set for the e-classes for abs and body:
            return cs-add-lib(f, CS[abs], CS[body])
    else:
        if there exists a non-empty cost set for each of a_0, ..., a_i in CS:
            let s_0, ..., s_i = CS[a_0, ..., a_i]
            return [(ls, c + 1) | ls, c in cs-product(s_0, ..., s_i)]
    return $\{\}$

# Returns at most M library selections.
# Removes all library selections containing more than $n$ libraries.
# $a$ is the input CostSet
# $M$ is the maximum size of a CostSet
# $n$ is the number of libraries to learn at a time
def cs-prune($a, M, n$):
    sort_by(a.libs, full-cost)
    for a-sel in $a$:
       if len(a.libs) >= n:
          a.remove(a-sel)
    retain($a$, M)
    return sort($a$)
\end{lstlisting}
\caption{Helper functions for extraction in \babble.}
\label{fig:helper3}
\end{figure}
